\documentclass[11pt]{article}
\usepackage[margin=1in]{geometry}
\usepackage{amsmath}
\usepackage{booktabs}
\usepackage{url}
\title{Synchronizing the manuscript and the documentation with d1c31d6}
\date{Prepared October 5, 2026}
\begin{document}
\maketitle

\section{Question}
The manuscript \path{paper/manuscript.tex} described the controller of
October 2 (\texttt{a29ccb6}): a cost-to-go CLF, a freely placed cruise core,
exact states, no road constraint, and a fourteen-run campaign with one
collision. Several documents described still older designs. The user asked for
the documentation and then the manuscript to be brought in line with the code
at \texttt{d1c31d6}.

\section{What was read}
The manuscript, every controller and configuration source, the estimator
runtime and its theory documents, and the reports
\path{JOINT_DESIGN_20261005.tex} and
\path{ESTIMATOR_DIRECT_TARGET_20261005.tex} were read line by line. The
perception sources, \path{scripts/}, \path{tests/} and the other reports of
October 1 to 5 were surveyed by three read-only assistants whose summaries
were used for the documentation changes; statements taken from them were
checked against the code where they enter the manuscript.

\section{Manuscript changes}
\begin{itemize}
\item Abstract, end of the introduction, system overview, architecture
  caption: the two components are evaluated together; the controller consumes
  the enclosure and imposes road edges; perception is not in the loop.
\item Obstacle perception: rectangle pose fit on lower-body points, started
  from the estimator's predicted pose.
\item Estimator: new subsection on the force-balance lateral velocity
  (monotone force balance, certified interval, the changed velocity row) and
  a description of the published ego zonotope and target parameter set.
  Lemmas, theorem and gain selection are unchanged.
\item Control design: rewritten. Variable horizon, shifted-plan and
  potential-field anchors, distance-dual rows with a least-penetration axis at
  overlap, road rows, rear-adhesion and sideslip cones, separating and
  road-recoverable terminal set with one proposition, quadratic CLF with its
  LMI synthesis, per-hold use of the enclosure, the two programs, the
  innovation-based trust scale, and the scope of the guarantees.
\item Results: Sections VI-A and VI-B (estimator alone) are kept and marked
  as measured before the model-aided measurement, whose kinematic branch they
  exercise. The closed-loop protocol, outcomes, in-loop estimation, the two
  runs without a solution, the window-fit ablation and the computation times
  are new. Discussion and conclusion are rewritten.
\item Two references were added (Zhai et al.\ 2024; Beal and Gerdes 2013).
\end{itemize}

\section{Sources of the numbers}
No simulation, estimator trial or unit test was run for this revision.
\begin{itemize}
\item Outcomes, gaps, road margins and recovery times:
  \path{ESTIMATOR_DIRECT_TARGET_20261005/campaign-metrics.csv}.
\item Failure table: \path{ESTIMATOR_DIRECT_TARGET_20261005/failures.csv} and
  the text of that report. Horizon statistics: \path{horizon-steps.csv}.
  Audit radii: \path{audit.csv}.
\item Positive primary optima, fresh anchors, input and state extremes,
  horizon of the exact runs, call-time shares and CLF decrease counts:
  derived from the recorded traces under
  \path{~/.cache/collisionAvoidance/estimator-direct-20261005} by
  \path{MANUSCRIPT_ALGORITHM_SYNC_20261005/deriveClosedLoopStatistics.py}
  (output \path{closed-loop-statistics.json}). The CLF counts quoted in the
  manuscript restrict the pairs to $V>10^{-3}$ and were computed with the same
  trace fields: 3064 pairs, 2520 with zero modeled slack, 73 of those without
  the requested decrease, none increasing.
\item Sequential timing (9.2~ms median, 41~ms 95th percentile, 3.4\% over
  50~ms) and the comparison with the kinematic measurement:
  \path{JOINT_DESIGN_20261005.tex}.
\item CLF contraction factors: \path{config/clfMatrices.json}.
\item Figure \path{closed_loop_results_v02}:
  \path{MANUSCRIPT_ALGORITHM_SYNC_20261005/figure-sources/closed_loop_results_v02.py}.
\end{itemize}

\section{Documentation changes}
\path{README.md} (rewritten; the previous text is kept as
\path{report/ROOT_README_SUPERSEDED_20261005.md}), \path{CLAUDE.md} and the
project-structure section of \path{AGENTS.md}, \path{perception/README.md},
\path{scripts/README.md}, \path{report/README.md} (current-state block and an
index of October 2 to 5), \path{paper/README.md},
\path{controller/PCBF_CLF_ARCHITECTURE.md} and
\path{controller/CONTROLLER_FILES.md} (multipliers at overlap),
\path{estimator/OBSERVER_ISS_THEORY.md} (links to removed files, pointer to
the model-aided measurement), and the comments of
\path{config/estimatorControllerIntegrationConfig.m}. No executable statement
was changed.

One disagreement between prose and code was found beyond stale names:
\path{PCBF_CLF_ARCHITECTURE.md} stated that an overlapping anchor keeps a zero
multiplier, while \texttt{solvePredictiveControl} calls
\texttt{supportLinearization}, which substitutes the feasible dual of the
least-penetrated rectangle axis. The document and the manuscript now describe
the code.

\section{Open items for review}
\begin{itemize}
\item The new derivations in the manuscript (force-balance interval and
  velocity row, collision tightening, CLF synthesis) restate the repository
  documents and code; they have not been independently reviewed.
\item The monotonicity of each axle force in yaw rate, speed, stiffness and
  friction, on which the corner evaluation of the force interval rests, is
  asserted as in \path{estimator/MODEL_AIDED_ESTIMATION.md}; the manuscript
  gives no proof.
\item The bibliographic data of the two added references should be checked
  against the publishers' records; the first names of the authors of Zhai et
  al.\ are given as initials.
\item The estimator-alone comparison was not rerun at \texttt{d1c31d6}.
\item Computation times were not re-measured at \texttt{d1c31d6}; the
  campaign timings were recorded with eight runs sharing the machine.
\item The architecture drawing still labels the road-boundary branch ``not
  yet used''; this remains true for perceived curbs only, and the caption says
  so.
\item \path{scripts/auditJointPredictiveSafety.py} checks a 4-m road
  half-width that no longer matches the MATLAB road; it was documented, not
  changed.
\end{itemize}
\end{document}
